% a06_WN_a: part 1 of 4 of the walled reading of Ba3 Prop. 1 (input from sec/a06_BALtorus_d7b.tex): the walled model and its zero-field kernels.
% Continues in sec/a06_WN_b.tex (classes, the two families, the statement), sec/a06_WN_c.tex (the phi phi A node clauses) and sec/a06_WN_d.tex (the bound and the constants).

\paragraph{The walled model.} Fix a step $s$, $1\le s\le K$, and work in scale-$s$ units: $\eta=L^{-s}$, $t_s\defeq L^s\eps$, $e_s\defeq e(L^s\eps)$ and $\lambda_s\defeq\lambda(L^s\eps)$ the running couplings of \cite[(1.33)]{Ba3}, and $T_\eta$ the torus, aligned with the blocks, with all sides $\ge1$. Let $v$ be a real function on the sites of $T_\eta$ with
\[
0\le v(x)\le\bar v\qquad\text{for every }x,
\]
and nothing else assumed about $v$: no smoothness, no constancy on any set, no symmetry. The walled scalar operator at zero background is
\[
H^v\defeq-\Delta^\eta_{T_\eta}+v+a_sP_s,\qquad G^v_s(T_\eta,0)\defeq(H^v)^{-1},
\]
with $a_sP_s$ the block term of the step-$s$ scalar propagator of \cite{Ba1}; at the constant $v=m^2t_s^2$ this is the bulk line $G_s(T_\eta,0)$ of Step~3(D). The wall of the rotation part is $v=(m^2+n\psi(t_s\,\cdot))t_s^2$ with a profile $0\le\psi\le1$, so that $\bar v=(m^2+n)t_s^2\le m^2+n$ once $t_s\le1$; the one-parameter family $v_\theta=(m^2+\theta n\psi(t_s\,\cdot))t_s^2$, $\theta\in[0,1]$, which interpolates between the bulk and the wall, satisfies the same bound. The \emph{walled model} is the perturbative scheme of \cite[Sect.~1]{Ba3} in which every scalar line $G_s(T_\eta,0)$, every slice \cite[(2.6)]{Ba3} and every zero-field object built from them is replaced by its $v$-analogue, while the vector lines are unchanged (the torus vector propagator carries no wall), the background is zero, and the counterterms are the model's own: a divergent two-scalar-leg subgraph $\gamma$ is renormalised by the vertex \cite[(1.7)]{Ba3} with the coefficient $\delta m^2_{s,\gamma}[v](x)$ given by Ba{\l}aban's rule \cite[p.~417]{Ba3} applied to the walled graph ($\gamma$ with both external legs at $x$ and its other leg point summed over the torus, with the same lines and the same labels, \cref{conv:a06d:LOC}, (Agr)), and the vacuum counterterm is $E_s[v]$, i.e.\ \cite[(1.24)]{Ba3} with walled lines. These counterterms depend on $x$, exactly as Ba{\l}aban's own do (\cite[p.~418]{Ba3}: the propagators ``are not translation invariant, so we have the dependence on $x$''), and the rule never uses translation invariance. Finite \cite[(1.7)]{Ba3} vertices with bounded $x$-dependent coefficients $c_v(x)$ are admitted as well; in the rotation part these are the non-renormalising parts of the inner mass counterterms and the \emph{wall vertex}, a two-scalar-leg vertex at $z$ with coefficient $-n\psi(t_sz)t_s^2$, which arises when a line of the bulk model is differentiated with respect to the wall height. The lines of a graph of the walled model are the full lines $G^v_s(T_\eta,0)$; the slices, orderings and regrouped lines of Step~3(D$_4$) are internal to the bounds.

\paragraph{Kernel notation.} For a kernel $K(x_1,x_2)$ on a lattice of spacing $\xi\in(0,1]$ and $S,T\subset\{1,2\}$, $D_S$ applies one forward difference $\partial^\xi_\mu$ in the variable $x_i$ for each $i\in S$, and $\Delta_T$ applies one H\"older quotient $|x_i-\tilde x_i|^{-\alpha}[K(\dots,x_i,\dots)-K(\dots,\tilde x_i,\dots)]$ for each $i\in T$ (in a fixed order; $\alpha\in(0,1)$ is the H\"older exponent of Step~3(D)); $\rho\defeq|x_1-x_2|$. Put $h_p(r)\defeq r^{-p}$ for $p>0$, $h_0(r)\defeq1+|\log r|$ and $h_p(r)\defeq1$ for $p<0$. In this notation the bounds of \cref{lem:a06d:S47}(b) on a remainder kernel $R$ read
\begin{equation}\label{eq:a06:WNR}
|\Delta_TD_SR(x_1,x_2)|\le C\,h_{|S|+\alpha|T|-1}(\xi+\rho)\,e^{-\delta\rho}\prod_{i\in T}|x_i-\tilde x_i|^\alpha,\qquad|S|,|T|\le2,
\end{equation}
and the zero-field, zero-mass form of the group-(A) estimates (\cite[Lemma~2.4, (2.35)--(2.36)]{Ba4}, \cref{thm:a06:A5}), extended to a difference in each variable by the same contour-shift argument, one factor $(e^{i\xi p_\mu}-1)/\xi$ per difference, for the pinned free operator $T\defeq-\Delta^\xi_X+bP_Y$ ($X$ a union of unit blocks with Neumann conditions, a periodic box or $\xi\Z^3$; $Y$ the unit blocks, or the $L$-blocks when $X$ is a union of $L$-blocks; $P_Y$ the block-average projection; $b\in[b_-,b_+]\subset(0,\infty)$) reads
\begin{equation}\label{eq:a06:WNZ1}
|\Delta_TD_ST^{-1}(x_1,x_2)|\le C_1(\xi+\rho)^{-p(S,T)}e^{-\delta\rho}\prod_{i\in T}|x_i-\tilde x_i|^\alpha,\qquad p(S,T)\defeq1+|S|+\alpha|T|,
\end{equation}
at every pair of points of $X$, with $T\ge\gamma'>0$ and $C_1,\delta,\gamma'$ depending only on $d=3$, $L$, $b_\pm$ and $\alpha$. (A quotient over two points at distance $>1$ is bounded by the two values, so \eqref{eq:a06:WNZ1} is needed only at $|x_i-\tilde x_i|\le1$; the same convention applies to every bound below.)

\begin{lemma}[lattice convolution]\label{lem:a06:WNK}
Let $0\le p_1,p_2<3$ and $\delta>0$. Then for all $x,y\in\xi\Z^3$ and $\xi\in(0,1]$,
\[
\sum_{z\in\xi\Z^3}\xi^3(\xi+|x-z|)^{-p_1}e^{-\delta|x-z|}(\xi+|z-y|)^{-p_2}e^{-\delta|z-y|}\ \le\ C_K\,h_{p_1+p_2-3}(\xi+|x-y|)\,e^{-\delta|x-y|/2},
\]
with $C_K=C_K(p_1,p_2,\delta)$. The same holds on a torus or a Neumann domain, the sum running over its sites.
\end{lemma}

\begin{proof}
Put $r\defeq|x-y|$. Since $|x-z|+|z-y|\ge r$, the exponentials are $\le e^{-\delta r/2}e^{-\delta(|x-z|+|z-y|)/2}$. On $\{|x-z|\le r/2\}$ one has $\xi+|z-y|\ge(\xi+r)/2$, and $\sum_{|x-z|\le r/2}\xi^3(\xi+|x-z|)^{-p_1}\le C(\xi+r)^{3-p_1}$ because $p_1<3$; the region $\{|z-y|\le r/2\}$ is symmetric; on the rest both factors are $\le C(\xi+|z-x|)^{-(p_1+p_2)}$ up to the exponentials, and $\sum_{|z|\ge r/2}\xi^3(\xi+|z|)^{-(p_1+p_2)}e^{-\delta|z|/2}$ is $\le C(\xi+r)^{3-p_1-p_2}$ if $p_1+p_2>3$, $\le C(1+|\log(\xi+r)|)$ if $p_1+p_2=3$ and $\le C$ otherwise. For $\xi+r\ge1$ the exponential factor dominates. On a torus or a Neumann domain the sites are a subset of $\xi\Z^3$ with the quotient (resp.\ restricted) distance, and the same sums bound the result.
\end{proof}

\begin{lemma}[walled zero-field kernels]\label{lem:a06:WN1}
Let $X$, $Y$, $b$ and $T=-\Delta^\xi_X+bP_Y$ be as in \eqref{eq:a06:WNZ1}, let $w\colon X\to[0,\bar v]$ and $T_w\defeq T+w$. Then:
\begin{enumerate}[label=\textup{(\alph*)},nosep,leftmargin=*]
\item\textup{(algebra)} $T_w\ge T\ge\gamma'$. The function $w$ commutes with every multiplication operator ($e^{\beta\rho}$ for a Lipschitz $\rho$, cutoffs, the gauge factors $V_\lambda$) and cancels in every difference of two operators carrying the same $w$; hence the gauge covariance $H^w(A^\lambda)=V_\lambda H^w(A)V_\lambda^*$, the commutator identities with cutoffs, the slice identity \cite[(2.6)]{Ba3} and its rescaling hold for the walled operator verbatim, the rescaled potential at scale $s_j$ being $s_j^2w(s_j\,\cdot)\in[0,\bar v s_j^2]\subset[0,\bar v]$, and the slices of the walled propagator are $G^{v,\eta}_{(j)}=a_j^2s_j^{-4}G^v_jQ_j^*C^{(j),v}Q_jG^v_j$ with $G^v_j\defeq(-\Delta+v+a_js_j^{-2}P_j)^{-1}$ and $C^{(j),v}$ the block operator built from $G^v_j$. Every walled zero-field scalar kernel is a scalar kernel times $\mathbf 1_N$, because $w$ is a scalar.
\item\textup{(the walled kernel bound)} $T_w^{-1}$ satisfies \eqref{eq:a06:WNZ1} at every pair of points, with $C_1$ replaced by $C_1^w\defeq C(d,L,\alpha,b_+,\gamma',\delta,C_1)(1+\bar v)^2$ and $\delta$ replaced by $\delta^w\defeq\min\{\delta/4,\beta'/\sqrt3\}/4$, where $\beta'\le1$ is the Combes--Thomas rate of \cref{lem:a07:CT} for $T_w$ (it depends only on $\gamma'$, $b_+$ and $L$, not on $w$).
\item\textup{(block sums)} For $|S|,|T|\le1$ and every block $B(y)$: $|\Delta_TD_S(T_w^{-1}\mathbf 1_{B(y)})(x)|\le c^we^{-\delta^w|x-y|}\prod_{i\in T}|x-\tilde x|^\alpha$.
\item\textup{(the remainder)} On $\xi\Z^3$ or a periodic box let $C^\xi\defeq(-\Delta^\xi+1)^{-1}$ be the free unit-mass kernel of Step~3(D$_4$) and $R^w\defeq T_w^{-1}-C^\xi$. Then $R^w$ satisfies \eqref{eq:a06:WNR} at every pair of points with $C=C_R(1+\bar v)^3$ and $\delta=\delta_R$, $C_R$ and $\delta_R$ depending only on $d,L,\alpha,b_\pm,\gamma'$ and the constants of \eqref{eq:a06:WNZ1}. In particular $R^w$ is a remainder line in the sense of \cref{lem:a06d:S47}(b), with constant $C_R(1+\bar v)^3$, and so are its slices and the kernels of \cref{lem:a06d:S47}(c) built from it.
\item\textup{(slices, cube kernels)} The slice bounds \cite[(2.10)--(2.12)]{Ba3} and the kernel bounds of the cube step of Step~3(D) hold for the walled slices $G^{v,\eta}_{(j)}$, including the two-ended entries and the mixed second difference, with constants depending polynomially on $\bar v$ and otherwise as in the bulk; the Combes--Thomas bound of \cref{lem:a07:CT} holds for $H^v$ with the $v$-free constants.
\item\textup{(periodisation)} $G^v_s(T_\eta,0)(x,x')=\sum_{n}G^{v_{\mathrm{per}}}_s(\eta\Z^3,0)(\tilde x,\tilde x'+n)$, the sum over the period lattice of $T_\eta$, $\tilde x,\tilde x'$ lifts of $x,x'$ and $v_{\mathrm{per}}$ the periodic extension of $v$; the sum converges absolutely, and the same identity holds slice by slice.
\end{enumerate}
\end{lemma}

\begin{proof}
(a) $w$ is a real scalar multiplication operator: it is self-adjoint and $\ge0$, it commutes with every other multiplication operator, and it cancels in $H^w(A)-H^w(A')$ and in $[H^w,\chi]$. The slice identity is the resolvent algebra of $-\Delta+w$ with the pin terms, which uses no property of the zeroth-order operator beyond coercivity; rescaling by $s_j$ multiplies a potential by $s_j^2$ and composes it with the dilation. Since $v$ acts as $v\mathbf 1_N$, the inverse of $H^v$ at zero field is a scalar kernel times $\mathbf 1_N$.

(b) \emph{Step 1 (Combes--Thomas).} $T_w=D^*D+w+bP_Y$ is of the form of \cref{lem:a07:CT} with $V=w\ge0$, $\Pi_B$ the block-average projections and $T_w\ge\gamma'$; so $\norm{e^{\beta\rho}T_w^{-1}e^{-\beta\rho}}\le2/\gamma'$ for every $1$-Lipschitz $\rho$ and $0\le\beta\le\beta'$, $\rho$ being $\sqrt3$-Lipschitz for the sup-metric, so that $\beta\le\beta'/\sqrt3$ is imposed, with $\beta'$ depending only on $\gamma'$, $b_+$ and the block size.
\emph{Step 2 (identity).} $T_w^{-1}=T^{-1}-T^{-1}wT^{-1}+T^{-1}wT_w^{-1}wT^{-1}$, which combines $T_w^{-1}=T^{-1}-T^{-1}wT_w^{-1}$ with $T_w^{-1}=T^{-1}-T_w^{-1}wT^{-1}$.
\emph{Step 3 (no operations).} Put $r\defeq|x-y|$. The first term is $\le C_1(\xi+r)^{-1}e^{-\delta r}$ by \eqref{eq:a06:WNZ1}. The second is $\le\bar vC_1^2C_Kh_{-1}(\xi+r)e^{-\delta r/2}=\bar vC_1^2C_Ke^{-\delta r/2}$ by \cref{lem:a06:WNK} with $p_1=p_2=1$. For the third write $a(z)\defeq T^{-1}(x,z)w(z)$, $b(z')\defeq w(z')T^{-1}(z',y)$ and $\rho\defeq|\cdot-x|$; then $|\ip a{T_w^{-1}b}|\le\norm{e^{\beta\rho}a}_2\,(2/\gamma')\,\norm{e^{-\beta\rho}b}_2$, and since $|z'-x|\ge r-|z'-y|$, $\norm{e^{-\beta\rho}b}_2\le e^{-\beta r}\norm{e^{\beta|\cdot-y|}b}_2$; both weighted norms are $\le\bar vC_1J$ with $J^2\defeq\sup_\xi\sum_z\xi^3(\xi+|z|)^{-2}e^{-(\delta-2\beta)|z|}<\infty$ (here $d=3$ and $\beta\le\delta/4$). Hence $|T_w^{-1}(x,y)|\le C^{(1)}(1+\bar v)^2(\xi+r)^{-1}e^{-\delta'r}$ with $\delta'\defeq\min\{\delta/2,\beta\}$ (for $\xi+r\le1$ the constant terms are $\le(\xi+r)^{-1}$, for $\xi+r\ge1$ the exponential absorbs the power).
\emph{Step 4 (general $S,T$).} Apply $\Delta_TD_S$ to the identity of step 2. The operations at $x_1$ act only on the first factor $T^{-1}(x_1,\cdot)$ and those at $x_2$ only on the last factor $T^{-1}(\cdot,x_2)$; $w$ and the middle factor $T_w^{-1}$ are never differentiated. Put $p_i\defeq1+[i\in S]+\alpha[i\in T]\in[1,2+\alpha]$, so that $p(S,T)=p_1+p_2-1$. The first term is \eqref{eq:a06:WNZ1}. The second is, by \eqref{eq:a06:WNZ1} for the two factors and \cref{lem:a06:WNK}, $\le2^{|T|}\bar vC_1^2C_Kh_{p_1+p_2-3}(\xi+\rho)e^{-\delta\rho/2}\prod|x_i-\tilde x_i|^\alpha$. For the third insert the bound of step 3 for the middle factor and apply \cref{lem:a06:WNK} twice: the convolution of $(\xi+|x_1-z|)^{-p_1}$ with $(\xi+|z-z'|)^{-1}$ gives $h_{p_1-2}(\xi+|x_1-z'|)$, which is $\le C_\alpha(\xi+|x_1-z'|)^{-q}$ with $q\defeq\max\{p_1-2,0\}+\alpha/2<3$ (this covers the logarithm at $p_1=2$), and the second convolution with $(\xi+|z'-x_2|)^{-p_2}$ gives $h_{q+p_2-3}(\xi+\rho)$. Since $p_1+p_2-3<p(S,T)$ and $q+p_2-3<p(S,T)$ with $p(S,T)\ge1$, both $h$-factors are $\le C(\xi+\rho)^{-p(S,T)}$ for $\xi+\rho\le1$, and for $\xi+\rho\ge1$ the exponentials dominate. Each application of \cref{lem:a06:WNK} halves the rate: the second term decays at the rate $\delta/2$ and the third at $\delta'/4$. Taking $\beta\defeq\min\{\delta/4,\beta'/\sqrt3\}$, so that $\delta'=\beta$, and collecting the constants gives $C_1^w=C(1+\bar v)^2$ and the rate $\delta^w=\min\{\delta/2,\beta/4\}=\beta/4$.

(c) Sum (b) over $x'\in B(y)$: since $p(S,T)\le2+\alpha<3$, $\sum_{x'\in B(y)}\xi^3(\xi+|x-x'|)^{-p}e^{-\delta^w|x-x'|}\le C(p)e^{-\delta^w\dist(x,B(y))}$.

(d) Since $(C^\xi)^{-1}=-\Delta^\xi+1$ and $T_w^{-1}$ exists, $R^w=T_w^{-1}\big((C^\xi)^{-1}-T_w\big)C^\xi=T_w^{-1}(1-w)C^\xi-b\,T_w^{-1}P_YC^\xi$. This is the identity \cite[p.~437]{Ba3} ($G^\xi(0)=C^\xi+G^\xi(0)(1-m^2-aP)C^\xi$) with the constant $1-m^2$ replaced by the function $1-w(z)$, which is never differentiated: the operations at $x_1$ fall on $T_w^{-1}(x_1,\cdot)$, those at $x_2$ on $C^\xi(\cdot,x_2)$. In the first term, (b) for $T_w^{-1}$, the free bound \eqref{eq:a06:WNZ1} for $C^\xi$ (it is the case $b=0$, $m^2=1$ of the group-(A) estimates), $|1-w|\le1+\bar v$ and \cref{lem:a06:WNK} give $C(1+\bar v)^3h_{p_1+p_2-3}(\xi+\rho)e^{-\delta_R\rho}\prod|x_i-\tilde x_i|^\alpha$ with $p_1+p_2-3=|S|+\alpha|T|-1$. The second term is $b\sum_Y|Y|^{-1}(T_w^{-1}\mathbf 1_Y)(x_1)(\mathbf 1_Y^{\mathsf T}C^\xi)(x_2)$, and (c) for $T_w^{-1}\mathbf 1_Y$ together with the block-sum bound for $C^\xi\mathbf 1_Y$ give $Ce^{-\delta_R\rho}\prod|x_i-\tilde x_i|^\alpha$ after the sum over $Y$, which is $\le h_{|S|+\alpha|T|-1}(\xi+\rho)$ for $\xi+\rho\le1$. The slices of $R^w$ and the kernels of \cref{lem:a06d:S47}(c) are obtained from these bounds exactly as in \cref{lem:a06d:S47}, which uses only the bounds \eqref{eq:a06:WNR}.

(e) The slice $G^{v,\eta}_{(j)}$ is the composition of (a) with $G^v_j$ bounded by (b) at scale $j$ (unit-block units after rescaling, with the potential $s_j^2v(s_j\,\cdot)\in[0,\bar v]$) and with the block-operator kernel $C^{(j),v}$, whose exponential bound uses only the coercivity $H^v\ge H^0\ge\gamma_0$ and the Combes--Thomas argument of \cref{lem:a07:CT}, so its constants are $v$-free; the bulk proof of \cite[(2.10)--(2.12)]{Ba3} in Step~3(D) is a composition of these three bounds and of lattice sums, and the mixed second difference and the two-ended entries are the cases $|S|=2$ of (b). The kernel bounds of the cube step are sums of slice bounds over the scales. The Combes--Thomas bound is \cref{lem:a07:CT} with $V=v$.

(f) The right side is well defined because the image sum converges absolutely by (b) on $\eta\Z^3$ with the periodic potential $v_{\mathrm{per}}$ (which lies in $[0,\bar v]$); it is periodic, and it solves the torus equation because the $\eta\Z^3$ operator with $v_{\mathrm{per}}$ commutes with the period translations; the torus inverse is unique. The slice identity (a) is compatible with the periodisation term by term.
\end{proof}
